\section*{अध्याय \ref{ch:b1:fluid-envelopes-structure} --- वायुमंडल और महासागर की संरचना}
% values: calc/ch26.py

\begin{solution}{exo:b1:fluid-envelopes-structure:1}
(क)~$H = R^*T/(M_0 g_0) = 8314.32 \times 250/(28.9644 \times 9.80665) =
\qty{7.32}{km}$। (ख)~$e^{10/7.32} = 3.92$: दाब लगभग एक-चौथाई तक गिरता है। (ग)~$H$ तापमान के अनुपात में है: \qty{216.65}{K} पर वह \qty{6.34}{km} है, जबकि \qty{288.15}{K} पर \qty{8.43}{km}। किसी दिए गए दाब पर ठंडी हवा सघन होती है, इसलिए दाब ऊँचाई के साथ तेज़ी से गिरता है।
\end{solution}

\begin{solution}{exo:b1:fluid-envelopes-structure:2}
$\Gamma_d = \qty{9.76}{K/km}$ के साथ: \qtyrange{0}{0.3}{km}, $11.0 > 9.76$, अस्थायी; \qtyrange{0.3}{1.2}{km}, $9.7$, आँकड़ों की सटीकता के भीतर उदासीन; \qtyrange{1.2}{1.5}{km}, एक व्युत्क्रमण; \qtyrange{1.5}{8}{km}, स्थायी; \qty{8}{km} से ऊपर, समतापीय, स्थायी। व्युत्क्रमण सबसे स्थायी है, क्योंकि वहाँ $\Gamma_d - \Gamma = 13.8$ सबसे बड़ा है।
\end{solution}

\begin{solution}{exo:b1:fluid-envelopes-structure:3}
(क)~$\sigma = 1026.95 - 1000 = \qty{26.95}{kg.m^{-3}}$। (ख)~B, जो हल्का है, A के ऊपर पड़ता है। (ग)~सभी समुद्री जल के घनत्व एक ही अंकों, 10, से शुरू होते हैं; $\sigma$ केवल वही अंक रखता है जो एक जल से दूसरे में बदलते हैं।
\end{solution}

\begin{solution}{exo:b1:fluid-envelopes-structure:4}
(क)~$S = 0.25 \times 35.40 + 0.75 \times 34.30 = 34.575$; $\theta = 0.25 \times
15.0 + 0.75 \times 4.0 = \qty{6.75}{\celsius}$। (ख)~खंड PQ पर, Q से P की ओर एक-चौथाई रास्ते पर (\cref{prop:b1:fluid-envelopes-structure:mixing})। (ग)~P, Q और R के मिश्रण त्रिभुज PQR को भरते हैं। $(34.60, 8.0) = a\,\mathrm P + b\,\mathrm Q + c\,\mathrm R$ लिखकर, जिसमें $a + b + c = 1$ है, $a = 0.345$, $b = 0.788$, $c = -0.133$ मिलते हैं: R का एक ऋणात्मक अनुपात, इसलिए बिंदु त्रिभुज के बाहर है और तीनों का कोई मिश्रण उसे नहीं बना सकता।
\end{solution}

\begin{solution}{exo:b1:fluid-envelopes-structure:5}
(क)~$\theta = 253.15 \times 2^{2/7} = \qty{308.6}{K}$। (ख)~\qty{1000}{hPa} पर, $\theta = T = \qty{298.15}{K}$। \omterm{def:b1:fluid-envelopes-structure:potential-temperature}{विभव तापमान} ऊपर की ओर बढ़ता है, \qty{298.2} से \qty{308.6}{K} तक: शुष्क हवा के लिए परत स्थायी है। (ग)~वह अपने \omterm{def:b1:fluid-envelopes-structure:potential-temperature}{विभव तापमान}, \qty{308.6}{K}, पर पहुँचती, जो \qty{35.4}{\celsius} है, वहाँ की हवा से दस अंश गरम।
\end{solution}

\begin{solution}{exo:b1:fluid-envelopes-structure:6}
(क)~$T = 288.15 - 6.5 \times 5 = \qty{255.65}{K}$ और $p = \num{101325} \times
(255.65/288.15)^{5.256} = \qty{54020}{Pa}$, \qty{540.2}{hPa}। (ख)~$H =
\qty{8.43}{km}$ और $p = \num{101325}\,e^{-5/8.43} = \qty{56010}{Pa}$। वास्तविक (मानक) दाब कम है: वहाँ की हवा \qty{288.15}{K} से ठंडी है, इसलिए सघन, और दाब ऊँचाई के साथ तेज़ी से गिरता है। (ग)~$p = \num{22632}\,
e^{-9/6.34} = \qty{5475}{Pa}$, \cref{tab:b1:fluid-envelopes-structure:layers} का मान।
\end{solution}

\begin{solution}{exo:b1:fluid-envelopes-structure:7}
(क)~$\beta\,\Delta S = \alpha\,\Delta T$ से $\Delta S = 1.67/7.61 = 0.22$। (ख)~$\Delta\rho = 1026.95 \times (-1.67\times10^{-4} \times 2 + 7.61\times10^{-4}
\times 0.3) = \qty{-0.11}{kg.m^{-3}}$: जल हल्का हो जाता है। (ग)~$\Delta\rho = 1028.11 \times (-0.53\times10^{-4} \times 2 + 7.85\times10^{-4}
\times 0.3) = \qty{+0.13}{kg.m^{-3}}$: वही परिवर्तन अब उसे सघन बनाते हैं, क्योंकि हिमांक के निकट जल में गरम होना घनत्व को मुश्किल से बदलता है।
\end{solution}

\begin{solution}{exo:b1:fluid-envelopes-structure:8}
(क)~$N^2 = (9.81/260) \times (9.76 - 4)\times10^{-3} = \qty{2.17e-4}{s^{-2}}$,
$N = \qty{0.0147}{s^{-1}}$, आवर्तकाल $2\pi/N = \qty{426}{s}$, लगभग सात मिनट। (ख)~$N^2 = (9.81/270) \times 14.76\times10^{-3} = \qty{5.36e-4}{s^{-2}}$,
$N = \qty{0.0232}{s^{-1}}$, आवर्तकाल \qty{271}{s}। (ग)~$N$ शून्य की ओर जाता है और आवर्तकाल अनंत तक बढ़ता है: उदासीन परत में प्रत्यानयन बल लुप्त हो जाता है।
\end{solution}

\begin{solution}{exo:b1:fluid-envelopes-structure:9}
(क)~यह सघन है, $\sigma_\theta = 27.20$ बनाम 26.30: AAIW अधिक ताज़ा है पर बहुत ठंडा। (ख)~$\Delta\sigma = 1026.95 \times [-1.67\times10^{-4} \times
(4.31 - 15.43) + 7.61\times10^{-4} \times (34.30 - 35.54)] = 1.907 - 0.969 =
0.94$, जबकि सटीक रूप से $0.90$। (ग)~तापमान: शीतलन \qty{1.91}{kg.m^{-3}} जोड़ता है, ताज़गी \qty{0.97}{kg.m^{-3}} घटाती है।
\end{solution}

\begin{solution}{exo:b1:fluid-envelopes-structure:10}
(क)~\cref{thm:b1:fluid-envelopes-structure:stability} से, $(g/\theta)\,\mathrm
d\theta/\mathrm dz = (g/T)(\Gamma_d - \Gamma)$; $\mathrm d\theta/\mathrm dz = 0$ से $\Gamma = \Gamma_d$। (ख)~$T = T_s - \Gamma_d z$, स्थिर प्रवणता $L = \Gamma_d$ की एक परत: \cref{prop:b1:fluid-envelopes-structure:layer-pressure} से, $p =
p_s (T/T_s)^{gM/(R\Gamma_d)}$, और $gM/(R\Gamma_d) = gMc_p/(Rg) = c_pM/R$
($= 3.5$ शुष्क हवा के लिए)। (ग)~$p$ वहाँ शून्य होता है जहाँ $T = 0$, $z = T_s/\Gamma_d =
300/9.76 = \qty{30.7}{km}$ पर। एकसमान \omterm{def:b1:fluid-envelopes-structure:potential-temperature}{विभव तापमान} की वास्तविक परतें अधिकतम कुछ किलोमीटर गहरी होती हैं: उनकी ऊपरी सीमा स्थायी हवा से होती है, जैसे क्षोभमंडल की समतापमंडल से।
\end{solution}

\begin{solution}{exo:b1:fluid-envelopes-structure:11}
(क)~$S = (33.787 + 36.050)/2 = 34.919$, $\theta = \qty{8.0}{\celsius}$। (ख)~रैखिक अवस्था-समीकरण के साथ $\sigma$, $S$ और $\theta$ का एक रैखिक फलन है, इसलिए मिश्रण का $\sigma$ दोनों का माध्य है, 27.00। (ग)~$\alpha$ \qty{2}{\celsius} पर \qty{0.75e-4}{K^{-1}} है और \qty{14}{\celsius} पर \qty{2.07e-4}{K^{-1}}, जबकि $\beta$ थोड़ा बदलता है (\num{7.80e-4} और \num{7.53e-4})। तापमान बढ़ने के साथ घनत्व अधिकाधिक तेज़ी से गिरता है: वह तापमान का एक अवतल फलन है। मिश्रण-रेखा के साथ, मिश्रण माध्य तापमान पर है, जहाँ घनत्व दोनों जनकों को जोड़ने वाली जीवा के ऊपर है, यानी 27.00 से ऊपर। (घ)~अधिकता वक्रता से आती है, एक द्वितीय-कोटि का पद: वह जनकों के तापमान अंतर के वर्ग के अनुपात में बढ़ती है।
\end{solution}

\begin{solution}{exo:b1:fluid-envelopes-structure:12}
(क)~दाब का अंतर $5604 - 5500 = \qty{104}{dbar}$ है: पिंड $0.129 \times 0.104 = \qty{0.013}{\celsius}$ ठंडा होता है, \qty{0.718}{\celsius} तक। (ख)~\qty{5400}{m} पर उसके चारों ओर का जल \qty{0.730}{\celsius} पर है, गरम: लगभग समान लवणता वाला पिंड परिवेश से सघन है और वापस नीचे डूबता है। परत स्थायी है, यद्यपि मापा गया तापमान नीचे की ओर बढ़ता है; उसका \omterm{def:b1:fluid-envelopes-structure:potential-temperature}{विभव तापमान} नीचे की ओर घटता है। (ग)~$\theta \approx
0.7313 - 0.086 \times 5.604 = \qty{0.249}{\celsius}$, सटीक मान के एक हज़ारवें अंश के भीतर।
\end{solution}

\begin{solution}{pb:b1:fluid-envelopes-structure:1}
\begin{enumerate}
  \item तापमान \qty{50}{m} तक लगभग एकसमान है, \qty{300}{m} तक तेज़ी से गिरता है, \qty{1000}{m} तक धीमे, और नीचे लगभग एकसमान है। लवणता सतह पर 35.80 से \qty{800}{m} पर न्यूनतम 34.32 तक घटती है, फिर \qty{3000}{m} पर 34.90 तक बढ़ती है।
  \item $22.40 - 0.5 = \qty{21.90}{\celsius}$ \qty{50}{m} (22.10) और \qty{75}{m} (19.20) के बीच पहुँचता है: $50 + 25 \times 0.20/2.90 = \qty{52}{m}$।
  \item $(22.10 - 13.00)/250 = \qty{0.0364}{\celsius/m}$, प्रति \qty{100}{m} \qty{3.6}{\celsius}। \omterm{def:b1:fluid-envelopes-structure:ocean-layers}{ताप-सीमा-परत} \omterm{def:b1:fluid-envelopes-structure:ocean-layers}{मिश्रित परत} के आधार से, \qty{50}{m} के निकट, लगभग \qty{800}{m} तक फैली है; वह \qty{50}{m} और \qty{75}{m} के बीच सबसे तीखी है (प्रति \qty{100}{m} \qty{11.6}{\celsius})।
  \item $\sigma = 1026.95 \times [1 - 1.67\times10^{-4} \times 12.4 +
        7.61\times10^{-4} \times 0.80] - 1000 = 25.45$, \qty{0.74}{kg.m^{-3}} अधिक ऊँचा। सतही जल संदर्भ अवस्था से बहुत दूर है, और $\alpha$ \qty{22}{\celsius} पर \qty{10}{\celsius} से कहीं बड़ा है: रैखिक रूप गर्मी से होने वाले हल्केपन को कम आँकता है।
  \item \qty{800}{m} पर लवणता न्यूनतम, \qty{4.8}{\celsius} पर 34.32, अंटार्कटिक मध्यवर्ती जल है, दक्षिणी महासागर की वर्षा-वाली पश्चिमी हवाओं के नीचे बना। \qty{3000}{m} पर अधिकतम, \qty{2.5}{\celsius} पर 34.90, उत्तर अटलांटिक गहन जल है, ग्रीनलैंड और लैब्राडोर सागरों में बना।
  \item $(2.50 - 2.26)/3.042 = \qty{0.079}{\celsius}$ प्रति \qty{1000}{dbar}।
  \item $\sigma_\theta$ गहराई के साथ बढ़ता है: स्तंभ स्थायी है। \omterm{def:b1:fluid-envelopes-structure:ocean-layers}{घनत्व-सीमा-परत} वहाँ है जहाँ $\sigma_\theta$ सबसे तेज़ बढ़ता है, ऊपरी कुछ सौ मीटर में, \omterm{def:b1:fluid-envelopes-structure:ocean-layers}{ताप-सीमा-परत} के साथ: \qty{200}{m} और \qty{800}{m} के बीच \num{0.82}, नीचे के \qty{2200}{m} में केवल \num{0.70}।
  \item $S = (35.45 + 34.32)/2 = 34.885$; $\theta = (14.90 + 4.80)/2 =
        \qty{9.85}{\celsius}$।
  \item C: 26.461; I: 27.310; मिश्रण: 26.886, ठीक दोनों का माध्य (26.886), क्योंकि रैखिक रूप रैखिक है।
  \item जनकों का माध्य $(26.342 + 27.159)/2 = 26.751$ है; मिश्रण, 26.888 पर, \qty{0.14}{kg.m^{-3}} सघन है। यह कैबेलिंग है।
  \item रैखिक रूप के साथ, परिच्छेद \qty{300}{m} पर 26.61, \qty{500}{m} पर 26.88 और \qty{800}{m} पर 27.31 देता है: मिश्रण का घनत्व \qty{500}{m} के निकट पड़ता है। वहाँ का जल, $(34.82, \qty{9.6}{\celsius})$, मिश्रण $(34.885, \qty{9.85}{\celsius})$ के बहुत निकट है: केंद्रीय जल और AAIW के बीच का परिच्छेद स्वयं उनके बीच की एक मिश्रण-रेखा के निकट है।
  \item $S = 0.25 \times 35.45 + 0.75 \times 34.32 = 34.60$, $\theta = 0.25
        \times 14.90 + 0.75 \times 4.80 = \qty{7.33}{\celsius}$। \qty{500}{m} और \qty{800}{m} के बीच परिच्छेद का प्रक्षेप \qty{650}{m} के निकट $(34.57, \qty{7.2}{\celsius})$ देता है: लगभग उसी पर।
  \item \si{K/km} में: A--B 10.0; B--C 9.7; C--D $-5.5$; D--E 6.29; E--F 6.50;
        F--G 6.57; G--H 0.33; H--I $-0.50$; I--J $-0.83$।
  \item A--B अस्थायी; B--C उदासीन; C--D एक व्युत्क्रमण; D--E, E--F, F--G स्थायी; G--H लगभग समतापीय, स्थायी; H--I और I--J व्युत्क्रमण। क्षोभसीमा G पर है, \qty{12.5}{km} और \qty{-59.0}{\celsius} पर, जहाँ तापमान घटना बंद करता है।
  \item औसत तापमान \qty{214.65}{K} है: $H = 287 \times 214.65/9.81 =
        \qty{6.28}{km}$, और $147\,e^{-4/6.28} = \qty{77.7}{hPa}$, मापे गए 78 से सुसंगत।
  \item A पर हवा, \qty{22.0}{\celsius}, समुद्र से ठंडी है: ऊष्मा समुद्र से हवा की ओर बहती है, जो नीचे से गरम होती है। परत A--B, जिसकी तापह्रास दर $\Gamma_d$ से ऊपर है, उसका परिणाम है।
  \item B पर: $22.0 - 9.76 \times 0.5 = \qty{17.1}{\celsius}$, 17.0 के मुकाबले; C पर \qty{7.4}{\celsius}, 7.3 के मुकाबले; D पर \qty{3.5}{\celsius}, 9.5 के मुकाबले। पिंड C तक हवा से थोड़ा गरम है, और D पर बहुत ठंडा।
  \item व्युत्क्रमण के आधार पर, \qty{1.5}{km} के निकट: व्युत्क्रमण में प्रवेश करते ही पिंड अपने चारों ओर की हवा से ठंडा, इसलिए सघन हो जाता है, जो उसका उठना रोकता है।
  \item अस्थायी सतही परत से मुक्त पिंड उसमें ऊपर उठते हैं और उसे तब तक मिलाते हैं जब तक उसका \omterm{def:b1:fluid-envelopes-structure:potential-temperature}{विभव तापमान} एकसमान न हो जाए, जिससे उसकी तापह्रास दर $\Gamma_d$ हो जाती है; ऊपर का व्युत्क्रमण मिश्रण रोकता है।
  \item K में: A 294.15; B 293.99; C 294.08; D 300.51; E 312.65; F 329.29; G
        346.28; H 369.50; I 446.99; J 600.91।
  \item $\theta$ A से B तक घटता है (अस्थायी), B से C तक लगभग स्थिर है (उदासीन) और ऊपर बढ़ता है (स्थायी), C से D तक और G के ऊपर सबसे तेज़: प्रश्न~14 के समान वर्गीकरण।
  \item E--F: $(329.29 - 312.65)/4 = \qty{4.2}{K/km}$; H--I: $(446.99 -
        369.50)/4 = \qty{19.4}{K/km}$, लगभग पाँच गुना अधिक।
  \item माध्य $\theta = \qty{408.2}{K}$ के साथ: $N^2 = (9.81/408.2) \times
        0.01937 = \qty{4.66e-4}{s^{-2}}$, $N = \qty{0.0216}{s^{-1}}$, आवर्तकाल \qty{291}{s}, लगभग पाँच मिनट।
  \item $N^2 = (9.81/1027) \times (26.35 - 24.79)/150 = \qty{9.9e-5}{s^{-2}}$,
        $N = \qty{0.0100}{s^{-1}}$, आवर्तकाल \qty{630}{s}। निचला समतापमंडल अधिक स्थायी है, $N^2$ \omterm{def:b1:fluid-envelopes-structure:ocean-layers}{ताप-सीमा-परत} की तुलना में लगभग पाँच गुना बड़ा।
  \item G पर, $\theta = 214.15 \times (1000/186)^{2/7} = \qty{346}{K}$: क्षोभसीमा पर \omterm{def:b1:fluid-envelopes-structure:potential-temperature}{विभव तापमान} लगभग \qty{346}{K} है।
\end{enumerate}
\end{solution}
